\documentclass[11pt]{article}
\usepackage{amsmath,amssymb}
\usepackage{graphicx}
\usepackage{hyperref}
\usepackage{geometry}
\geometry{margin=1in}

\title{VERALUME: A Meta-Control Framework for Regulating LLM Responses under Semantic Tension}
\author{Christian Duguay}
\date{}

\begin{document}
\maketitle

\begin{abstract}
Large Language Models (LLMs) often exhibit unstable behavior when confronted with prompts involving paradoxes, self-reference, or high semantic tension, leading to forced resolutions, hallucinations, or incoherent explanations. We introduce VERALUME, a lightweight meta-control framework designed to regulate LLM response behavior without modifying model parameters. VERALUME operates as an external control system that estimates semantic tension in user inputs, dynamically selects a response regime, and adapts control parameters through feedback based on measured response characteristics. We formalize VERALUME as a discrete-time dynamical system with continuous state variables representing mobility, elasticity, semantic tension, and response chiralité. This work makes no claims about internal cognition or consciousness; instead, it provides an analyzable and reproducible method for steering LLM behavior under epistemically unstable conditions.
\end{abstract}

\section{Introduction}
Large Language Models have demonstrated strong capabilities in reasoning and generation, yet they frequently fail under paradoxical or self-referential prompts. Rather than resolving such prompts, models often hallucinate coherence. This paper reframes the issue as a control problem: how to regulate response regimes under semantic tension.

\section{Problem Definition}
We treat the LLM as a black-box stochastic function:
\[
y = \text{LLM}(x \mid \pi, \tau)
\]
where $x$ is the prompt, $\pi$ a system instruction policy, and $\tau$ a temperature parameter.

\section{VERALUME Framework}
\subsection{State Space}
At interaction step $t$, the controller state is:
\[
S_t = (M_t, E_t, \Delta_t, \chi_t)
\]
with all variables in $[0,1]$.

\subsection{Semantic Tension Estimation}
Semantic tension is computed via a heuristic function:
\[
\Delta_t = \sigma(f(x_t) - \mu)
\]
where $f$ captures lexical and structural indicators of instability.

\subsection{Control Law}
A deformation force is defined as:
\[
F_t = \frac{E_t}{\max(M_t, \varepsilon)} \cdot \Delta_t
\]
Thresholds on $F_t$ determine the response regime.

\paragraph{Geometric Interpretation.} The scalar tension $\Delta_t$ acts as a curvature term on the semantic manifold, locally contracting the Fisher Information Metric $g_{ij}(\theta)$ defined as:
\[
g_{ij}(\theta) = \mathbb{E}\left[\left(\frac{\partial}{\partial\theta_i}\log p(x|\theta)\right)\left(\frac{\partial}{\partial\theta_j}\log p(x|\theta)\right)\right]
\]
where $\theta$ represents the latent state and $p(x|\theta)$ the token distribution. The force $F_t$ is the counter-acting vector required to maintain geodesic stability under semantic curvature.

\subsection{Response Generation}
\[
y_t = \text{LLM}(x_t \mid \pi(F_t), \tau(\Delta_t))
\]

\subsection{Chiralit\'e Measurement}
Response chiralit\'e is measured as:
\[
\chi_t = 1 - \text{sim}(x_t, y_t)
\]
where $\text{sim}$ denotes semantic similarity.

\subsection{Adaptive Update}
\[
M_{t+1} = \text{clip}(M_t + \alpha (1 - \chi_t))
\]
\[
E_{t+1} = \text{clip}(E_t + \beta \chi_t)
\]

\section{Cartographic Analysis}
\begin{figure}[htbp]
    \centering
    \includegraphics[width=0.85\textwidth]{state_space_placeholder.png}
    \caption{\textbf{Trajectoire dans l'Espace de Phase VERALUME ($M \times E$).} 
    Le graphique illustre une transition de phase induite par une injection de dissonance sémantique. 
    Le système débute dans un état cristallin (Bas-Gauche, Régime M1) caractérisé par une faible mobilité. 
    Sous l'effet de la tension $\Delta_{cog}$ (taille des points), l'élasticité $E_t$ augmente, permettant une bifurcation vers un état fluide (Haut-Droite, Régime M4). 
    Le gradient de couleur indique l'augmentation de la Chiralité $\chi_t$ (orthogonalité de la réponse).}
    \label{fig:veralume_phase_space}
\end{figure}

\section{Simulated Experimental Results}
Using simulated trajectories, VERALUME reveals model-specific behavioral signatures characterized by attractors and regime transitions.

\section{Discussion and Limitations}
VERALUME regulates observable behavior only and does not provide guarantees of correctness or claims about cognition.

\section{Conclusion}
We presented VERALUME, a meta-control framework that stabilizes LLM interaction under semantic tension through adaptive regime management.

\begin{thebibliography}{9}
\bibitem{brown2020} Brown et al., Language Models are Few-Shot Learners, NeurIPS 2020.
\bibitem{ouyang2022} Ouyang et al., Training language models to follow instructions, NeurIPS 2022.
\bibitem{astrom2008} \AA str\"om and Murray, Feedback Systems, Princeton 2008.
\end{thebibliography}

\end{document}
